% Source: https://artofproblemsolving.com/community/c3h1597766_geometric_probability_question
% Stable identity: aops-topic-1597766-q01
The midpoint of $(x_1,y_1)$ and $(x_2,y_2)$ is $\left(\dfrac{x_1+x_2}{2},\dfrac{y_1+y_2}{2}\right)$. Its coordinates are integers exactly when $x_1\equiv x_2$ and $y_1\equiv y_2\pmod 2$, and then they automatically lie between $1$ and $3$, so the midpoint is in $S$.

Group the points by the parities of their coordinates:
\begin{itemize}
\item (odd, odd): $(1,1),(1,3),(3,1),(3,3)$, four points, giving $\binom42=6$ pairs;
\item (odd, even): $(1,2),(3,2)$, giving $1$ pair; (even, odd): $(2,1),(2,3)$, giving $1$ pair;
\item (even, even): only $(2,2)$, giving no pair.
\end{itemize}

\begin{asy}
usepackage("amsmath");
size(4.5cm);
defaultpen(fontsize(10pt));
for(int i=1;i<=3;++i) for(int j=1;j<=3;++j){
  pen p=(i%2==1&&j%2==1)?red:((i%2==1)?blue:((j%2==1)?heavygreen:black));
  dot((i,j),p+linewidth(5));}
draw((1,1)--(3,3),red+dashed);draw((1,3)--(3,1),red+dashed);draw((1,1)--(1,3),red+dashed);draw((3,1)--(3,3),red+dashed);draw((1,1)--(3,1),red+dashed);draw((1,3)--(3,3),red+dashed);
draw((1,2)--(3,2),blue+dashed);draw((2,1)--(2,3),heavygreen+dashed);
label("$(1,1)$",(1,1),SW);label("$(3,3)$",(3,3),NE);
label("\mbox{}",(min(currentpicture,user=true)+max(currentpicture,user=true))/2,black); // keeps the text after the figure black when LaTeX typesets the labels inline
\end{asy}

(The dashed segments are the $8$ good pairs.) So there are $6+1+1=8$ good pairs out of $\binom92=36$ pairs in all, and the probability is
\[
\dfrac{8}{36}=\boxed{\dfrac{2}{9}}.
\]
The poster's count of $8$ good pairs was right; the error was the total: $\binom92=36$, not $45$ ($45=\binom{10}{2}$).
